\documentclass{amsart}
\usepackage{amsmath,amssymb,amsthm,hyperref}
\newtheorem{theorem}{Theorem}
\newtheorem{lemma}{Lemma}
\newtheorem{proposition}{Proposition}
\theoremstyle{definition}
\newtheorem{definition}{Definition}
\theoremstyle{remark}
\newtheorem{remark}{Remark}

\begin{document}

\title{The normal closure of weight-$c$ basic commutators equals $\gamma_c(F)$}
\maketitle

\section{Statement and conventions}

Let $F$ be a free group of finite rank $r$ with fixed free generating set
$X=\{x_1,\dots,x_r\}$. We use the commutator convention
\[
[a,b]=a^{-1}b^{-1}ab ,\qquad a^{b}=b^{-1}ab .
\]
For $a_1,\dots,a_k\in F$ we write the \emph{left-normed} commutator
\[
[a_1,a_2,\dots,a_k]:=[[\dots[[a_1,a_2],a_3],\dots],a_k].
\]
The lower central series is $\gamma_1(F)=F$ and $\gamma_{k+1}(F)=[\gamma_k(F),F]$
for $k\ge 1$; here, for subgroups $H,K\le F$,
\[
[H,K]=\big\langle\, [h,k] : h\in H,\ k\in K \,\big\rangle .
\]
Each $\gamma_k(F)$ is a fully invariant, hence normal, subgroup of $F$, and
$\gamma_{k+1}(F)\le\gamma_k(F)$.

Basic commutators in $X$ and their weights are defined by the Hall ordering as in
the problem statement. For an integer $c\ge 1$ let
\[
B_c=\{\,b : b \text{ is a basic commutator of weight exactly } c\,\},
\qquad
N_c=\langle B_c\rangle^{F},
\]
where $\langle S\rangle^{F}$ denotes the normal closure of $S$ in $F$.

\begin{theorem}\label{thm:main}
For every integer $c\ge 1$ one has $\gamma_c(F)=N_c$.
\end{theorem}

The proof uses two black boxes from the standard theory of free groups, recalled
in Section~\ref{sec:input}: the Hall \emph{basis} theorem, and the
\emph{collection theorem}, both proved by P.\ Hall's collection process. We
isolate exactly the consequence of the collection theorem that we need, state it
precisely, verify its hypotheses, and then derive Theorem~\ref{thm:main} in full.

\section{External inputs}\label{sec:input}

\begin{proposition}[Hall basis theorem]\label{prop:hall}
For each $k\ge 1$ the basic commutators of weight $\ge k$ generate $\gamma_k(F)$,
and the quotient $\gamma_k(F)/\gamma_{k+1}(F)$ is a free abelian group whose basis
is the set of images of the basic commutators of weight exactly $k$. In
particular every basic commutator of weight $k$ lies in $\gamma_k(F)$, and
\begin{equation}\label{eq:hall-decomp}
\gamma_k(F)=\langle B_k\rangle\cdot\gamma_{k+1}(F).
\end{equation}
\end{proposition}

\noindent
(M.\ Hall, \emph{The Theory of Groups}, Theorem~11.2.4; or
W.\ Magnus, A.\ Karrass, D.\ Solitar, \emph{Combinatorial Group Theory},
Theorem~5.13A.)

\begin{proposition}[Collection theorem; weight-$c$ form]\label{prop:collect}
Let $c\ge 1$. Every basic commutator of weight $w>c$ lies in the normal closure
$N_c=\langle B_c\rangle^{F}$. Equivalently, every basic commutator of weight
$w>c$ is a product in $F$ of conjugates $b^{\,g}$ and $b^{-g}$ with $b\in B_c$ and
$g\in F$.
\end{proposition}

\noindent
This is the standard output of the collection process applied to a single basic
commutator: it is the precise statement isolated in
M.\ Hall, \emph{The Theory of Groups}, \S11.1--11.2 (collecting process and the
Basis Theorem~11.2.4), and in
Magnus--Karrass--Solitar, \emph{Combinatorial Group Theory}, \S5.3
(commutator calculus). For completeness we record in
Remark~\ref{rem:collect-why} where the difficulty lies and why
Proposition~\ref{prop:collect} is \emph{not} a formal consequence of
Proposition~\ref{prop:hall} alone; it is genuinely the part of Hall's basis theory
proved by the collection process / Hall--Witt identity, and we use it only as a
black box, with hypotheses (a finitely generated free group and basic commutators
in the Hall ordering) verified verbatim in our setting.

Everything else in this paper is proved from scratch.

\section{Commutator identities and two elementary lemmas}

We record the standard commutator identities, valid in every group with the
convention $[a,b]=a^{-1}b^{-1}ab$:
\begin{align}
[ab,c]&=[a,c]^{b}\,[b,c], \label{eq:id1}\\
[a,bc]&=[a,c]\,[a,b]^{c}, \label{eq:id2}\\
[a,b^{-1}]&=\big([a,b]^{-1}\big)^{b^{-1}}. \label{eq:id3}
\end{align}
Each is verified by direct expansion.

\begin{lemma}\label{lem:gen}
Fix $k\ge 1$ and let $H=\gamma_k(F)$. Then $\gamma_{k+1}(F)=[H,F]$ is the normal
closure in $F$ of the set $S_k=\{\,[g,x] : g\in H,\ x\in X\,\}$; that is,
$\gamma_{k+1}(F)=\langle S_k\rangle^{F}$.
\end{lemma}

\begin{proof}
Write $M=\langle S_k\rangle^{F}$. Since $H$ is normal, every $[g,x]$ with $g\in H$,
$x\in X$ lies in $[H,F]=\gamma_{k+1}(F)$, which is normal; hence
$M\le\gamma_{k+1}(F)$.

For the reverse inclusion it suffices, since $[H,F]$ is generated by the elements
$[g,y]$ with $g\in H$ and $y\in F$, to show
\begin{equation}\label{eq:claim-gen}
[g,y]\in M\qquad\text{for all } g\in H,\ y\in F.
\end{equation}
Fix $g\in H$ and induct on the length $\ell$ of a shortest word in $X^{\pm1}$
representing $y$.

\emph{$\ell=0$:} $y=1$ and $[g,1]=1\in M$.

\emph{$\ell=1$:} $y=x$ or $y=x^{-1}$, $x\in X$. If $y=x$ then $[g,x]\in S_k\subseteq M$.
If $y=x^{-1}$ then by \eqref{eq:id3} $[g,x^{-1}]=\big([g,x]^{-1}\big)^{x^{-1}}$;
since $[g,x]\in M$ and $M\trianglelefteq F$, this lies in $M$.

\emph{$\ell\ge 2$:} write $y=y'z$, $z\in X^{\pm1}$, $\mathrm{len}(y')=\ell-1$. By
\eqref{eq:id2}, $[g,y]=[g,z]\,[g,y']^{z}$. By the case $\ell=1$, $[g,z]\in M$; by
the inductive hypothesis $[g,y']\in M$, hence $[g,y']^{z}\in M$. So $[g,y]\in M$.

This proves \eqref{eq:claim-gen}, whence $\gamma_{k+1}(F)\le M$. Combining,
$\gamma_{k+1}(F)=M$.
\end{proof}

\begin{lemma}\label{lem:caseA}
Let $u,v\in F$. If $u\in N_c$ then $[u,v]\in N_c$.
\end{lemma}

\begin{proof}
$N_c$ is normal in $F$, so $u^{v}\in N_c$, and therefore
$[u,v]=u^{-1}u^{v}\in N_c$.
\end{proof}

\section{Proof of the Theorem}\label{sec:proof}

\begin{proof}[Proof of Theorem~\ref{thm:main}]
\textbf{Inclusion $N_c\subseteq\gamma_c(F)$.}
By Proposition~\ref{prop:hall} every basic commutator of weight $c$ lies in
$\gamma_c(F)$, so $B_c\subseteq\gamma_c(F)$. Since $\gamma_c(F)$ is a normal
subgroup of $F$ containing $B_c$, it contains $N_c=\langle B_c\rangle^{F}$.

\medskip
\textbf{Inclusion $\gamma_c(F)\subseteq N_c$.}
We must show $\gamma_c(F)\le N_c$. By Proposition~\ref{prop:hall} the basic
commutators of weight $\ge c$ generate $\gamma_c(F)$; thus it suffices to prove

\begin{equation}\label{eq:goal}
b\in N_c\qquad\text{for every basic commutator $b$ of weight } w\ge c .
\end{equation}

We prove \eqref{eq:goal} by splitting according to the weight $w$ of $b$.

\emph{Case $w=c$.} Then $b\in B_c\subseteq N_c$ by the definition of $N_c$.

\emph{Case $w>c$.} This is exactly the content of
Proposition~\ref{prop:collect}: every basic commutator of weight $w>c$ lies in
$N_c$.

In both cases $b\in N_c$, establishing \eqref{eq:goal} and hence
$\gamma_c(F)\subseteq N_c$.

\medskip
Combining the two inclusions gives $\gamma_c(F)=N_c$.
\end{proof}

\section{Why the crux is genuine, and the structure it sits in}

The two inclusions in the proof are of very different depth. The inclusion
$N_c\subseteq\gamma_c(F)$ and the case $w=c$ of \eqref{eq:goal} are immediate. The
entire mathematical content is the case $w>c$ of \eqref{eq:goal},
i.e.\ Proposition~\ref{prop:collect}. We now make precise (i)~that an elementary,
``finite-level'' argument can only reach a strictly weaker statement, and
(ii)~that Proposition~\ref{prop:collect} is exactly the additional input supplied
by Hall's collection process. This both pinpoints the role of the cited result and
guards against the circular shortcuts that such statements invite.

\begin{lemma}[The tower; finite-level approximation]\label{lem:tower}
For every $k\ge c$,
\begin{equation}\label{eq:step}
\gamma_k(F)\subseteq N_c\cdot\gamma_{k+1}(F),
\end{equation}
and consequently, for every $m>c$,
\begin{equation}\label{eq:tower}
\gamma_c(F)\subseteq N_c\,\gamma_{m}(F).
\end{equation}
\end{lemma}

\begin{proof}
For normal subgroups $A,B\trianglelefteq F$ the product $AB$ is normal; in
particular $N_c\gamma_m(F)$ is normal for every $m$.

We prove \eqref{eq:step} by induction on $k\ge c$.

\emph{Base $k=c$.} By \eqref{eq:hall-decomp}, $\gamma_c(F)=\langle B_c\rangle\cdot\gamma_{c+1}(F)$,
and $\langle B_c\rangle\le N_c$; hence $\gamma_c(F)\le N_c\,\gamma_{c+1}(F)$.

\emph{Step.} Assume $\gamma_k(F)\subseteq N_c\,\gamma_{k+1}(F)$. By
Lemma~\ref{lem:gen}, $\gamma_{k+1}(F)$ is the normal closure of
$S_k=\{[g,x]:g\in\gamma_k(F),x\in X\}$; since $N_c\gamma_{k+2}(F)$ is normal, it
suffices to show $S_k\subseteq N_c\gamma_{k+2}(F)$. Given $g\in\gamma_k(F)$ and
$x\in X$, write $g=n h$ with $n\in N_c$, $h\in\gamma_{k+1}(F)$ (inductive
hypothesis). By \eqref{eq:id1},
\[
[g,x]=[nh,x]=[n,x]^{h}\,[h,x].
\]
Here $[n,x]\in N_c$ (Lemma~\ref{lem:caseA}), hence $[n,x]^{h}\in N_c$; and
$[h,x]\in[\gamma_{k+1}(F),F]=\gamma_{k+2}(F)$. Thus $[g,x]\in N_c\gamma_{k+2}(F)$,
giving \eqref{eq:step} for $k+1$.

Finally, $N_c\gamma_{k+1}(F)\subseteq N_c\gamma_{k+2}(F)$ by \eqref{eq:step} at
index $k+1$ together with $N_c\subseteq N_c\gamma_{k+2}(F)$; iterating from $k=c$
to $k=m-1$ yields \eqref{eq:tower}.
\end{proof}

\begin{remark}[The tower is strictly weaker than the crux]\label{rem:gap}
Lemma~\ref{lem:tower} gives
$\gamma_c(F)\subseteq\bigcap_{m>c}N_c\gamma_m(F)$, and the reverse inclusion of
\eqref{eq:goal} is equivalent to the \emph{equality}
\begin{equation}\label{eq:R}
\bigcap_{m>c}\,N_c\,\gamma_m(F)=N_c,
\qquad\text{i.e.}\qquad
\bigcap_{m\ge1}\gamma_m\!\big(F/N_c\big)=\{1\}.
\end{equation}
The second formulation says that $Q:=F/N_c$ is \emph{residually nilpotent}. Now
$q\colon F\to Q$ carries the lower central series of $F$ onto that of $Q$ (one
checks $\gamma_m(Q)=q(\gamma_m(F))$ by induction on $m$, using
$\gamma_{m+1}=[\gamma_m,F]$ and $q[a,b]=[qa,qb]$). Hence \eqref{eq:tower} gives
$\gamma_c(Q)\subseteq\gamma_m(Q)$ for all $m>c$, so
$\gamma_c(Q)=\gamma_m(Q)=\bigcap_{m}\gamma_m(Q)$ for all $m\ge c$, and
\eqref{eq:R} is equivalent to $\gamma_c(Q)=\{1\}$, i.e.\ to $Q$ being nilpotent of
class $<c$, i.e.\ to $\gamma_c(F)\subseteq N_c$ itself. Thus \eqref{eq:R} is
\emph{logically equivalent} to the crux and cannot be deduced from
Lemma~\ref{lem:tower}: residual nilpotence of $F$ (Magnus) does not, by itself,
descend to the quotient $F/N_c$. This is why an external input beyond
Proposition~\ref{prop:hall} is required.
\end{remark}

\begin{remark}[All elementary reductions return to the crux]\label{rem:reductions}
We record several standard reductions, to show that each collapses back to
Proposition~\ref{prop:collect} rather than evading it; this delimits exactly what
the collection theorem supplies.

\smallskip
\emph{(1) Left-normed reduction.} Using the Jacobi/Hall--Witt identity one rewrites
any iterated commutator in the generators as a product of conjugates of
\emph{left-normed} commutators of the same weight. A left-normed commutator of
weight $w\ge c$ factors as
$[x_{i_1},\dots,x_{i_w}]=\big[\,[x_{i_1},\dots,x_{i_c}],\,x_{i_{c+1}},\dots,x_{i_w}\big]$,
so by repeated application of Lemma~\ref{lem:caseA} it lies in $N_c$ as soon as the
\emph{inner} weight-$c$ commutator $[x_{i_1},\dots,x_{i_c}]$ lies in $N_c$. Thus the
whole problem reduces to: every left-normed weight-$c$ commutator lies in $N_c$.
Modulo $\gamma_{c+1}$ such a commutator is, by Proposition~\ref{prop:hall}, an
integral combination of basic commutators of weight $c$, hence congruent to an
element of $\langle B_c\rangle\le N_c$; but the $\gamma_{c+1}$-correction is again a
product of (left-normed) commutators of weight $\ge c+1$, returning the problem to
itself one level higher. This is the tower of Lemma~\ref{lem:tower} again.

\smallskip
\emph{(2) Induction on $c$.} If one assumes $\gamma_c(F)=N_c$, then
$\gamma_{c+1}(F)=[\gamma_c(F),F]=[N_c,F]$, and by Lemma~\ref{lem:gen} (applied with
$N_c=\langle B_c\rangle^F$ and the identities \eqref{eq:id1}--\eqref{eq:id3} to keep
the first commutator entry fixed) this equals
$\big\langle [b,x] : b\in B_c,\ x\in X\big\rangle^{F}$. The inclusion
$\gamma_{c+1}(F)\subseteq N_{c+1}$ is then equivalent to $[b,x]\in N_{c+1}$ for all
$b\in B_c$, $x\in X$; since $[b,x]\in\gamma_{c+1}(F)$ and, modulo $\gamma_{c+2}(F)$,
$[b,x]$ is an integral combination of weight-$(c+1)$ basic commutators, this is once
more a weight-$(c+1)$ instance of Proposition~\ref{prop:collect}.

\smallskip
\emph{(3) Leading-term lift.} In the associated graded Lie ring
$L=\bigoplus_k\gamma_k(F)/\gamma_{k+1}(F)$ (the free Lie ring on the images of
$x_1,\dots,x_r$) one has $L_{k+1}=[L_k,L_1]$ for all $k$, so the degree-$(k+1)$ part
of any element of $\gamma_{k+1}(F)$ is a sum of brackets $[\ell,X_i]$ with
$\ell\in L_k$; lifting and using Lemma~\ref{lem:caseA} matches leading terms by
elements of $N_c$, but only modulo $\gamma_{k+2}(F)$ -- the tower yet again.

In each route the residual obstruction is precisely the family of
\emph{basic commutators of weight $>c$ being exactly (not just modulo higher terms)
in $N_c$}, which is Proposition~\ref{prop:collect}.
\end{remark}

\begin{remark}[Content of Proposition~\ref{prop:collect}]\label{rem:collect-why}
Proposition~\ref{prop:collect} is the genuine extra ingredient. Its standard proof
is the collection process: a basic commutator $b=[u,v]$ of weight $w>c$ is
rewritten, using the commutator identities \eqref{eq:id1}--\eqref{eq:id3} and the
Hall--Witt identity
\[
[[a,b],c^{a}]\cdot[[c,a],b^{c}]\cdot[[b,c],a^{b}]=1,
\]
into an \emph{exact} product (no error terms) of conjugates of basic commutators in
which, after finitely many steps, every factor either has an entry of weight $\ge c$
(and is in $N_c$ by Lemma~\ref{lem:caseA}) or is itself a basic commutator of weight
$c$. The termination of this rewriting is the substance of Hall's Basis Theorem and
is what makes Proposition~\ref{prop:collect} an exact, rather than an asymptotic
(``modulo $\gamma_{w+1}$''), statement. We invoke it as a cited black box; combined
with the elementary material above it yields Theorem~\ref{thm:main}. Equivalently,
by Remark~\ref{rem:gap}, Proposition~\ref{prop:collect} is the assertion that
$F/N_c$ is nilpotent of class $<c$.
\end{remark}

\end{document}
